%----------HEADING----------%
% 아이콘 미사용. FontAwesome 글리프는 텍스트 추출 시 깨진 문자(# §)로
% 나오므로, 각 항목에 일반 텍스트 라벨을 사용합니다.
\begin{center}
    % \textls 미사용. kotex가 \pickup@font를 재정의하기 때문에 microtype의
    % 자간 조절이 추출 시 글자 사이에 공백을 넣습니다 ("M i n j u n").
    % 영문 템플릿에서는 정상 동작하지만 여기서는 이름이 검색되지 않게 됩니다.
    \textbf{\LARGE \scshape \textcolor{accent}{김민준 (Minjun Kim)}} \\ \vspace{2pt}
    \vspace{5pt}
    \small
    전화: 010-1234-5678 $|$
    \href{mailto:minjun.kim@example.com}{이메일: \underline{minjun.kim@example.com}} $|$
    \href{https://github.com/minjunkim}{GitHub: \underline{github.com/minjunkim}} $|$
    서울특별시
\end{center}
